\chapter{概率图模型}
\label{chap:graphical-models}
\begin{quote}\small
\textit{本章摘要。}概率图模型用图表示随机变量之间的条件独立关系，用局部分布或势函数构造联合分布。本章从朴素贝叶斯的条件独立假设出发，把已知类别改为隐变量，建立高斯混合模型，再把单个隐变量扩展为状态序列，介绍状态空间模型及其推断。随后从局部归一化的最大熵马尔可夫模型转入全局归一化的条件随机场，解释无向图的分离、团分解和能量表达，最后介绍玻尔兹曼机与受限玻尔兹曼机。各模型的差异落实为边的方向、节点取值、条件独立关系、隐变量安排和分布参数化五个方面。
\end{quote}

温度和振动可能共同受到设备状态影响，相邻时刻的状态也可能相关。若直接列出所有变量的联合概率，离散变量的配置数会随变量个数指数增长。概率图模型（Probabilistic Graphical Model，PGM）通过明确哪些依赖需要保留，把联合分布拆成较小的部分\cite{bishop2006pattern,murphy2022pml}。图给出分布的结构，具体概率还需要由局部分布及其参数确定。

本章以朴素贝叶斯作为简单起点。先把类别视为已知标签，再把它变为需要推断的隐变量，随后让隐状态沿时间连接。接着改变建模目标，直接学习给定观测后的标签分布，并比较局部归一化与全局归一化。最后在无向图中引入二元隐变量，研究怎样通过限制连边获得较容易计算的模型。这是一条逐步改变假设的教学路线，各模型之间的联系将在相应分布上说明。

阅读时要区分随机变量、变量的观测值与模型参数。以下用大写$X,Y,Z$表示随机变量，小写$x,y,z$表示取值，$\theta$表示参数。$x_{1:T}$表示从第一个到第$T$个时刻的观测序列。推断是在参数给定时计算未知变量的分布，学习是用数据估计参数。贝叶斯公式是推断的基本工具，但图模型也可以用最大似然学习参数。

\section{朴素贝叶斯与有向图}
\subsection{类别如何把多个观测联系起来}
假设根据温度是否偏高$X_1$和振动是否偏高$X_2$，判断设备状态$Y\in\{0,1\}$，其中$0$为正常，$1$为退化。朴素贝叶斯（Naive Bayes，NB）假设：给定类别后，各观测特征条件独立。对$d$个特征，这一假设写成
\begin{equation}
 p(x_1,\ldots,x_d\mid y)=\prod_{j=1}^d p(x_j\mid y),
 \qquad p(y,x_1,\ldots,x_d)=p(y)\prod_{j=1}^d p(x_j\mid y).
 \label{eq:bayes-naive-joint}
\end{equation}
这里$p(y)$是类别先验，$p(x_j\mid y)$是类别给定时第$j$个特征的分布。二元特征可以使用伯努利分布，连续特征可以使用高斯分布，同一个图也可以同时包含这两类节点。

\begin{figure}[htbp]
\centering
\begin{tikzpicture}[>=Latex, node distance=1cm and 1.4cm,
 state/.style={circle,draw,minimum size=9mm},every node/.style={font=\small}]
 \node[state] (y) {$Y$};
 \node[state,fill=blue!12,below left=of y] (x1) {$X_1$};
 \node[state,fill=blue!12,below=of y] (x2) {$X_2$};
 \node[state,fill=blue!12,below right=of y] (xd) {$X_d$};
 \draw[->,thick] (y)--(x1); \draw[->,thick] (y)--(x2); \draw[->,thick] (y)--(xd);
\end{tikzpicture}
\caption{朴素贝叶斯的有向图。类别是各特征的共同父节点，给定类别后特征条件独立。训练时类别通常已知，预测时需要根据特征推断类别。}
\label{fig:bayes-naive-structure}
\end{figure}

图\ref{fig:bayes-naive-structure}中没有特征间的直接连边。它允许特征通过类别产生边缘相关。例如，设$p(Y=1)=0.2$，退化时温度和振动各以$0.8$概率偏高，正常时各以$0.1$概率偏高，并采用条件独立假设。则
\begin{align}
 p(X_1=1,X_2=1)&=0.2\times0.8^2+0.8\times0.1^2=0.136,\\
 p(X_1=1)p(X_2=1)&=(0.2\times0.8+0.8\times0.1)^2=0.0576.
\end{align}
两个数不同，说明温度和振动在不知道状态时相关。条件独立只在所给条件下成立。

观察到两个特征都偏高后，贝叶斯公式给出
\begin{equation}
 p(Y=1\mid X_1=1,X_2=1)
 =\frac{0.2\times0.8^2}{0.2\times0.8^2+0.8\times0.1^2}
 =\frac{16}{17}.
\end{equation}
分子衡量退化类别对观测的解释，分母对所有类别求和。若给定状态后两个传感器仍有共享噪声，该分解会重复计算相关证据，可能产生过于集中的类别概率。分类效果和概率校准都需要在独立数据上检验。

\subsection{有向无环图怎样定义分布}
朴素贝叶斯是贝叶斯网络（Bayesian Network，BN）的一种简单结构。贝叶斯网络使用有向无环图（Directed Acyclic Graph，DAG）。有向边$X_i\to X_j$把$X_i$列入$X_j$的父节点，无环表示沿箭头无法回到出发节点。

按拓扑顺序排列$m$个变量，使父节点先于子节点出现。概率链式法则给出
\begin{equation}
 p(x_1,\ldots,x_m)=\prod_{j=1}^m p(x_j\mid x_1,\ldots,x_{j-1}).
\end{equation}
再假设给定父节点后，$X_j$与其余先前变量条件独立，就得到贝叶斯网络的分解
\begin{equation}
 p(x_1,\ldots,x_m)=\prod_{j=1}^m p(x_j\mid x_{\operatorname{pa}(j)}).
 \label{eq:bayes-dag-factorization}
\end{equation}
$\operatorname{pa}(j)$表示节点$j$的父节点集合，没有父节点时使用边缘分布。每个因子都对自己的节点归一化，因而整个联合分布归一化。删去条件依赖是建模假设，需要数据或领域知识支持。箭头规定分布如何分解，因果解释还需要额外假设。

\begin{figure}[htbp]
\centering
\begin{tikzpicture}[>=Latex,state/.style={circle,draw,minimum size=9mm},every node/.style={font=\small}]
 \node[state] (a) at (0,0) {$A$};
 \node[state] (b) at (2,1) {$B$};
 \node[state] (c) at (2,-1) {$C$};
 \node[state] (d) at (4,0) {$D$};
 \draw[->,thick] (a)--(b); \draw[->,thick] (a)--(c);
 \draw[->,thick] (b)--(d); \draw[->,thick] (c)--(d);
 \node[anchor=west,align=left,text width=6cm] at (5,0) {拓扑顺序：$A,B,C,D$\\[4pt]
 父节点：$\operatorname{pa}(D)=\{B,C\}$\\[4pt]
 $p(a,b,c,d)$\\$=p(a)p(b\mid a)p(c\mid a)p(d\mid b,c)$};
\end{tikzpicture}
\caption{从有向图读取分解。每个节点贡献一个给定父节点的因子；$B,C$共同作为$D$的条件。}
\label{fig:bayes-dag-factorization}
\end{figure}

图\ref{fig:bayes-dag-factorization}给出一个具体例子。链式法则原本含$p(c\mid a,b)$与$p(d\mid a,b,c)$，图结构分别将它们简化为$p(c\mid a)$与$p(d\mid b,c)$。这两处简化对应$C\perp B\mid A$和$D\perp A\mid(B,C)$。读图时先列出每个节点的父节点，再将对应条件分布相乘。

\subsection{用d-分离判断条件独立}
条件独立$X\perp Y\mid S$表示：给定$S=s$后，知道$X$不会改变$Y$的条件分布。离散变量在$p(s)>0$时满足
\begin{equation}
 p(x,y\mid s)=p(x\mid s)p(y\mid s).
\end{equation}
连续变量用条件密度表达相同关系。条件集合$S$可以包含多个变量。仅看两点之间是否有边，无法判断它们是否条件独立，还需要检查连接它们的路径。

d-分离（d-separation）提供从DAG读出条件独立关系的规则。检查路径时，先忽略边的方向以寻找路径，再根据箭头在每个中间节点的会合方式判断路径是否阻断。三种基本结构如下。
\begin{enumerate}
 \item 链$X\to Z\to Y$或$X\leftarrow Z\leftarrow Y$。中间节点$Z$在条件集合中时，路径被阻断。例如知道当前状态后，更早状态不会经这条链额外影响下一状态。
 \item 分叉$X\leftarrow Z\to Y$。共同父节点$Z$在条件集合中时，路径被阻断。朴素贝叶斯的各特征属于这种情形。
 \item 碰撞$X\to Z\leftarrow Y$。只有当$Z$自身或它的某个后代在条件集合中时，这个中间节点才允许路径通过。后代指从$Z$沿箭头能够到达的节点。
\end{enumerate}

\begin{figure}[htbp]
\centering
\begin{tikzpicture}[>=Latex,state/.style={circle,draw,minimum size=8mm},every node/.style={font=\small}]
 \foreach \r/\title in {0/链,1/分叉,2/碰撞}{
  \node[state] (x\r) at (0,-1.8*\r) {$X$};
  \node[state] (z\r) at (2,-1.8*\r) {$Z$};
  \node[state] (y\r) at (4,-1.8*\r) {$Y$};
  \node[anchor=east] at (-0.6,-1.8*\r) {\title};
 }
 \draw[->,thick] (x0)--(z0); \draw[->,thick] (z0)--(y0);
 \draw[->,thick] (z1)--(x1); \draw[->,thick] (z1)--(y1);
 \draw[->,thick] (x2)--(z2); \draw[->,thick] (y2)--(z2);
 \node[anchor=west,align=left] at (5,0) {不观察$Z$：活跃\\观察$Z$：阻断};
 \node[anchor=west,align=left] at (5,-1.8) {不观察$Z$：活跃\\观察$Z$：阻断};
 \node[anchor=west,align=left] at (5,-3.6) {不观察$Z$及其后代：阻断\\观察$Z$或其后代：活跃};
\end{tikzpicture}
\caption{d-分离的三种基本路径结构。这里仅判断所画路径，假定端点不在条件集合中。存在其他路径时，还需逐条检查。}
\label{fig:bayes-d-separation}
\end{figure}

图\ref{fig:bayes-d-separation}显示，条件化对碰撞结构与另外两种结构的作用不同。前两类的中间节点称为非碰撞节点。给定条件集合$S$，一条路径活跃需要同时满足：路径上的所有非碰撞中间节点都不在$S$中，每个碰撞中间节点自身或至少一个后代在$S$中。若两个互不相交节点集合$A,B$之间的所有路径都被阻断，就称$A$和$B$被$S$ d-分离。$S$还应与$A,B$不相交。

实际判断时，可以先标记$S$及其祖先，再逐条检查$A$到$B$的路径。某个非碰撞节点属于$S$，或者某个碰撞节点不属于$S$且没有后代属于$S$，都会阻断该路径。所有路径阻断才能推出$X_A\perp X_B\mid X_S$。只找到一条阻断路径并不足够。

\paragraph{为什么观察碰撞节点会引入依赖。}
令$X,Y$为两个独立二元原因，各以$1/2$概率发生，$Z=1$表示至少一个原因发生。未观察$Z$时，$X$与$Y$独立。观察$Z=1$后，$(X,Y)$只可能是$(1,0),(0,1),(1,1)$且等概率，所以$p(Y=1\mid Z=1)=2/3$。若又知道$X=0$，则$p(Y=1\mid Z=1,X=0)=1$。选择了共同结果以后，两个原因产生依赖。若只观察$Z$的一个有信息的后代，也可能产生这种依赖。

\paragraph{为什么必须检查所有路径。}
考虑$A\to B\to D$和$A\to C\to D$组成的菱形DAG。给定$B$只阻断第一条路径，第二条仍活跃。给定$B,C$才阻断全部路径，从而得到$A\perp D\mid(B,C)$。d-分离保证该图所分解的所有分布都具有相应独立性。反过来，路径活跃只说明图结构不保证独立，某些特殊参数仍可能使变量独立。若要由分布中的独立性反推图结构，还需要忠实性等额外条件。


\section{高斯混合模型中的隐类别}
\subsection{把类别变成需要推断的变量}
若训练数据只有温度和振动，没有状态标签，就无法直接按类别统计$p(x\mid y)$。可以引入一个未观测的类别变量$Z\in\{1,\ldots,K\}$，用它表示样本属于哪个成分。隐变量（Latent Variable）是模型中的随机变量，未观测的取值需要根据数据推断。它与控制分布形状的参数不同。

若只把朴素贝叶斯的类别隐藏起来，得到的是隐类别模型。再规定每个类别下的连续观测服从高斯分布，才得到高斯混合模型（Gaussian Mixture Model，GMM）：
\begin{align}
 p(Z=k)&=\pi_k,\qquad \pi_k\geq0,\quad\sum_{k=1}^K\pi_k=1,\\
 p(x\mid Z=k)&=\mathcal N(x;\mu_k,\Sigma_k),\\
 p(x)&=\sum_{k=1}^K\pi_k\mathcal N(x;\mu_k,\Sigma_k).
 \label{eq:bayes-gmm-marginal}
\end{align}
其中$x,\mu_k\in\R^d$，$\Sigma_k\in\R^{d\times d}$为正定协方差矩阵。正定性保证每个成分具有非退化的多元高斯密度。生成一个样本时，先按$\pi$抽取成分，再从该成分的高斯分布抽取观测。每个样本有自己的$Z_i$，各样本在给定共享参数后独立。“单个隐变量”指每个样本只有一个成分指示变量。

若$\Sigma_k$为对角矩阵，给定成分后各特征独立，与高斯朴素贝叶斯具有相同的星形分解。一般GMM允许完整协方差，此时给定$Z$后观测分量仍可能相关。若把整个向量$X$画成一个节点，$Z\to X$不会表达向量内部是否独立，必须进一步说明协方差结构。引入隐变量本身不会强制高斯分布或对角协方差。

\begin{figure}[htbp]
\centering
\begin{tikzpicture}[>=Latex,state/.style={circle,draw,minimum size=10mm},every node/.style={font=\small}]
 \foreach \i in {1,2,3}{
  \node[state] (z\i) at (3*\i,1.5) {$Z_{\i}$};
  \node[state,fill=blue!12] (x\i) at (3*\i,0) {$X_{\i}$};
  \draw[->,thick] (z\i)--(x\i);
 }
 \node[align=center] at (6,2.6) {$Z_i\sim\operatorname{Categorical}(\pi)$：离散成分};
 \node[align=center] at (6,-1.1) {$X_i\mid Z_i=k\sim\mathcal N(\mu_k,\Sigma_k)$：连续观测};
\end{tikzpicture}
\caption{GMM的图结构，以三个独立样本为例。白色为隐变量，蓝色为观测变量。各样本共享参数$\pi,\mu_k,\Sigma_k$，图中参数视为给定，样本之间没有转移边。}
\label{fig:bayes-gmm-structure}
\end{figure}

图\ref{fig:bayes-gmm-structure}中，每条$Z_i\to X_i$连接表示$p(z_i)p(x_i\mid z_i)$。观测密度$p(x_i)$需要对该样本的成分$Z_i$求和。后面建立序列模型时，将在这些成分节点之间加入状态转移边。

\subsection{后验责任度与EM学习}
观察到$x_i$后，成分后验为
\begin{equation}
 r_{ik}=p(Z_i=k\mid x_i)
 =\frac{\pi_k\mathcal N(x_i;\mu_k,\Sigma_k)}
 {\sum_{\ell=1}^K\pi_\ell\mathcal N(x_i;\mu_\ell,\Sigma_\ell)}.
 \label{eq:bayes-gmm-responsibility}
\end{equation}
$r_{ik}$称为责任度，表示第$k$个成分对样本$i$的后验权重，各成分责任度之和为一。例如两个成分等权，均值分别为$0,2$，方差均为$1$。对$x=1$，两成分密度相同，责任度各为$1/2$。对$x=2$，第二个成分的责任度为$1/(1+e^{-2})\approx0.881$。这是保留成分不确定性的软分配。

训练要最大化观测对数似然
\begin{equation}
 \ell(\theta)=\sum_{i=1}^n\log\sum_{k=1}^K\pi_k\mathcal N(x_i;\mu_k,\Sigma_k).
\end{equation}
求和处于对数内部，使直接求解较困难。期望最大化（Expectation Maximization，EM）先用旧参数计算责任度，再以责任度为权重更新参数。令$N_k=\sum_i r_{ik}$，则
\begin{align}
 \pi_k^{\rm new}&=N_k/n,\qquad
 \mu_k^{\rm new}=\frac{1}{N_k}\sum_i r_{ik}x_i,\\
 \Sigma_k^{\rm new}&=\frac{1}{N_k}\sum_i r_{ik}
 (x_i-\mu_k^{\rm new})(x_i-\mu_k^{\rm new})^{\mathsf T}.
\end{align}
这里$x_i,\mu_k\in\R^d$，因此协方差更新的结果是$d\times d$矩阵。$N_k$是期望样本数，可以不是整数。更新对应最大化期望完整数据对数似然。精确E步和M步保证观测似然不下降，但不能保证全局最优。若$N_k=0$，更新无定义。若协方差收缩到零，未约束的高斯混合似然可能无界，因此需要限制协方差、设置下界或使用先验。成分编号可以任意置换，不能自动解释为正常、退化或故障。

\section{状态空间模型与隐状态序列}
\subsection{状态空间模型的结构假设}
GMM没有规定相邻样本的先后关系。若设备状态会持续一段时间，需要为每个时刻设置$Z_t$，再用状态转移连接它们。状态空间模型（State-Space Model，SSM）以隐状态序列解释观测序列。常用的一阶形式为
\begin{align}
 p(z_t\mid z_{1:t-1})&=p(z_t\mid z_{t-1}),\\
 p(x_t\mid z_{1:T},x_{1:t-1})&=p(x_t\mid z_t),\\
 p(z_{1:T},x_{1:T})&=p(z_1)p(x_1\mid z_1)
 \prod_{t=2}^T p(z_t\mid z_{t-1})p(x_t\mid z_t).
 \label{eq:bayes-ssm-joint}
\end{align}
第一式是一阶马尔可夫假设，当前状态只需通过上一状态读取过去。第二式是观测条件独立假设，给定状态序列后，各时刻观测独立，每个观测只依赖当前状态。它允许同一时刻的多个传感器采用联合观测分布。

齐次性是另一个假设：转移核$p(z_t\mid z_{t-1})$不随$t$改变。它不由一阶马尔可夫性推出，也不保证状态边缘分布不随时间变化。传感器关系不变时，还可假设观测分布参数在时间上共享。时变模型仍可以保留相同的链式图，改变的是局部分布参数。

隐状态可以离散，也可以连续。有限离散状态的模型通常称为隐马尔可夫模型。连续状态若满足线性高斯条件，可以用卡尔曼滤波精确推断。非线性或非高斯模型可以考虑粒子滤波近似推断。隐马尔可夫模型是一类模型，卡尔曼滤波和粒子滤波是推断算法，不能仅按名称把它们并列为三种图结构。

\begin{figure}[htbp]
\centering
\begin{tikzpicture}[>=Latex,state/.style={circle,draw,minimum size=12mm},every node/.style={font=\small}]
 \node[state] (z1) at (0,0) {$Z_{t-1}$};
 \node[state] (z2) at (4,0) {$Z_t$};
 \node[state] (z3) at (8,0) {$Z_{t+1}$};
 \node[state,fill=blue!12] (x1) at (0,-1.7) {$X_{t-1}$};
 \node[state,fill=blue!12] (x2) at (4,-1.7) {$X_t$};
 \node[state,fill=blue!12] (x3) at (8,-1.7) {$X_{t+1}$};
 \draw[->,thick] (z1)--node[above] {状态转移}(z2);
 \draw[->,thick] (z2)--node[above] {状态转移}(z3);
 \draw[->,thick] (z1)--(x1); \draw[->,thick] (z2)--(x2); \draw[->,thick] (z3)--(x3);
 \node[anchor=west] at (4.8,-0.85) {$p(x_t\mid z_t)$};
\end{tikzpicture}
\caption{状态空间模型的局部链结构。白色状态节点可离散或连续，蓝色节点为观测。线性高斯模型和HMM具有相同的依赖结构，区别在于节点取值和局部分布。}
\label{fig:bayes-ssm-structure}
\end{figure}

图\ref{fig:bayes-ssm-structure}中的状态连边是相对于独立GMM新增的依赖。图不规定转移是否齐次，也不规定局部分布是否高斯，这些条件需要另行说明。

\subsection{链结构怎样支持递归推断}
过滤（Filtering）计算$p(z_t\mid x_{1:t})$，只使用截至当前的观测。递推分为预测和更新：
\begin{align}
 p(z_t\mid x_{1:t-1})
 &=\int p(z_t\mid z_{t-1})p(z_{t-1}\mid x_{1:t-1})\,dz_{t-1},\\
 p(z_t\mid x_{1:t})
 &=\frac{p(x_t\mid z_t)p(z_t\mid x_{1:t-1})}
 {\int p(x_t\mid u)p(u\mid x_{1:t-1})\,du}.
 \label{eq:bayes-filter-recursion}
\end{align}
第一步将上一时刻的状态不确定性通过转移传播到当前，第二步用新观测重新加权。离散状态把积分改为求和。分母是预测观测密度，必须正且有限。不同过滤算法的区别，在于怎样表示状态分布并完成这两步运算。

\subsection{线性高斯模型与卡尔曼滤波}
设连续状态$z_t\in\R^q$，观测$x_t\in\R^d$，线性高斯模型为
\begin{equation}
 z_t=F_tz_{t-1}+\eta_t,\qquad x_t=H_tz_t+\varepsilon_t,
 \quad \eta_t\sim\mathcal N(0,Q_t),\quad
 \varepsilon_t\sim\mathcal N(0,R_t).
\end{equation}
$F_t\in\R^{q\times q}$描述状态转移，$H_t\in\R^{d\times q}$把状态映射为观测，$Q_t,R_t$分别描述过程噪声和观测噪声。假设初始状态高斯，各时刻噪声彼此独立且与初始状态独立，$R_t$正定。常数矩阵对应时间齐次的常用情形，时变矩阵也可使用以下递推。

若上一过滤分布为$\mathcal N(m_{t-1},P_{t-1})$，线性变换和独立高斯相加给出预测均值$m_t^-=F_tm_{t-1}$与协方差$P_t^-=F_tP_{t-1}F_t^{\mathsf T}+Q_t$。给定过去观测，状态与当前观测共同服从
\begin{equation}
 \begin{bmatrix}z_t\\x_t\end{bmatrix}\Bigm|x_{1:t-1}
 \sim\mathcal N\left(
 \begin{bmatrix}m_t^-\\H_tm_t^-\end{bmatrix},
 \begin{bmatrix}P_t^-&P_t^-H_t^{\mathsf T}\\
 H_tP_t^-&H_tP_t^-H_t^{\mathsf T}+R_t\end{bmatrix}\right).
\end{equation}
右下块是观测预测协方差，非对角块是状态与观测的交叉协方差。对这个联合高斯分布以当前观测$x_t$条件化，得到卡尔曼滤波（Kalman Filter，KF）\cite{kalman1960filter}：
\begin{align}
 m_t^-&=F_tm_{t-1},& P_t^-&=F_tP_{t-1}F_t^{\mathsf T}+Q_t,\\
 S_t&=H_tP_t^-H_t^{\mathsf T}+R_t,&
 K_t&=P_t^-H_t^{\mathsf T}S_t^{-1},\\
 m_t&=m_t^-+K_t(x_t-H_tm_t^-),&
 P_t&=P_t^--K_tS_tK_t^{\mathsf T}.
\end{align}
上标$-$表示接收当前观测之前，$S_t$为预测观测误差的协方差，$K_t$为卡尔曼增益。均值修正量由观测残差和两类不确定性的相对大小决定。实际计算用线性方程求解代替显式求逆，并使用稳定的协方差更新。

一维例子中，若预测状态均值为$0$、方差为$4$，观测$x_t=3$，$H_t=1$，观测噪声方差为$1$，则$K_t=4/5$，更新均值为$2.4$，方差为$0.8$。较精确的观测使更新更靠近观测。这里高斯分布在预测和更新后保持封闭，因此均值和协方差足以表示精确过滤分布。只有马尔可夫性与观测独立性，还不足以得到这个高斯递推。

\begin{figure}[htbp]
\centering
\begin{tikzpicture}[>=Latex,every node/.style={font=\small},box/.style={draw,rounded corners,align=center,minimum height=12mm,text width=2.7cm}]
 \node[box] (old) at (0,0) {上一过滤分布\\$\mathcal N(m_{t-1},P_{t-1})$};
 \node[box,fill=orange!12] (prior) at (4.2,0) {预测分布\\$\mathcal N(m_t^-,P_t^-)$};
 \node[box,fill=green!12] (post) at (8.4,0) {当前过滤分布\\$\mathcal N(m_t,P_t)$};
 \node[box,fill=blue!12] (obs) at (8.4,1.9) {新观测$x_t$\\观测矩阵$H_t$、噪声$R_t$};
 \draw[->,thick] (old)--node[above] {$F_t,Q_t$}(prior);
 \draw[->,thick] (prior)--node[above] {高斯条件化}(post);
 \draw[->,thick] (obs)--(post);
\end{tikzpicture}
\caption{卡尔曼滤波的计算流程。状态空间图仍如图\ref{fig:bayes-ssm-structure}，这里的方框表示分布与计算步骤，箭头表示计算次序。线性高斯条件使预测与更新后的分布始终保持高斯形式。}
\label{fig:bayes-kalman-recursion}
\end{figure}

图\ref{fig:bayes-kalman-recursion}把前面的递推分成两步。预测传播上一分布的均值与协方差，更新使用观测残差修正均值，并减少条件不确定性。

\subsection{非线性、非高斯模型与粒子滤波}
若状态转移或观测关系非线性，或者过滤分布有多个峰，仅传播均值和协方差可能遗漏重要结构。粒子滤波（Particle Filter，PF）用带权样本近似过滤分布：
\begin{equation}
 p(z_t\mid x_{1:t})\approx\sum_{s=1}^M w_t^{(s)}\delta_{z_t^{(s)}}(z_t),
 \qquad w_t^{(s)}\geq0,\quad\sum_s w_t^{(s)}=1.
\end{equation}
$M$为粒子数，$\delta_a$表示集中于$a$的单位点质量，每个粒子代表一种可能状态。这个表达是离散测度近似，不要求真实状态离散。

以转移分布为提议的自举粒子滤波，先从$p(z_t\mid z_{t-1}^{(s)})$传播每个粒子，再更新
\begin{equation}
 \widetilde w_t^{(s)}=w_{t-1}^{(s)}p(x_t\mid z_t^{(s)}),\qquad
 w_t^{(s)}=\frac{\widetilde w_t^{(s)}}{\sum_r\widetilde w_t^{(r)}}.
\end{equation}
这是式\eqref{eq:bayes-filter-recursion}的采样实现。如果先对祖先按旧权重重采样，旧权重已包含在祖先选择中，本步权重只需正比于观测似然。重采样复制高权重粒子并舍弃低权重粒子，之后权重设为$1/M$。它缓解权重集中，但也可能减少粒子多样性，不能凭空找回已遗漏的状态区域。

\begin{figure}[htbp]
\centering
\begin{tikzpicture}[>=Latex,every node/.style={font=\small}]
 \foreach \s/\title in {0/传播后的粒子,1/观测加权,2/重采样后}{
  \draw[gray] (4*\s,0)--(4*\s+2.6,0);
  \node at (4*\s+1.3,1.1) {\title};
 }
 \foreach \x in {0.2,0.8,1.7,2.3}{\fill[blue!65] (\x,0) circle (2pt);}
 \fill[blue!65] (4.2,0) circle (1.3pt);
 \fill[blue!65] (4.8,0) circle (2pt);
 \fill[blue!65] (5.7,0) circle (4.5pt);
 \fill[blue!65] (6.3,0) circle (1.3pt);
 \fill[blue!65] (8.8,0) circle (2pt);
 \foreach \y in {0,0.2,0.4}{\fill[blue!65] (9.7,\y) circle (2pt);}
 \draw[->,thick] (2.8,0.45)--(3.7,0.45);
 \draw[->,thick] (6.8,0.45)--(7.7,0.45);
 \node at (1.3,-0.65) {可能的状态位置};
 \node at (5.3,-0.65) {圆大小表示权重};
 \node at (9.3,-0.65) {复制后权重相等};
\end{tikzpicture}
\caption{粒子滤波的一维状态示意。观测似然改变权重；图示一次可能的重采样结果，高权重粒子被复制，低权重粒子可能消失。右侧同位置的三个粒子竖向错开以便显示，实际状态相同。重采样不是每步必须执行的操作。}
\label{fig:bayes-particle-filter}
\end{figure}

图\ref{fig:bayes-particle-filter}用圆的位置表示状态，用大小表示权重。它展示的是分布近似与计算过程，所依据的随机变量图仍是图\ref{fig:bayes-ssm-structure}。

有效粒子数常估计为$1/\sum_s(w_t^{(s)})^2$，可用于决定何时重采样。粒子数有限会带来采样误差，高维状态下权重退化可能很严重。粒子滤波不要求齐次转移，适用于能够采样转移、评价观测似然并满足相应支持条件的模型。


\subsection{有限离散状态与隐马尔可夫模型}
隐马尔可夫模型（Hidden Markov Model，HMM）把状态空间限制为有限集合。其前向后向、最优路径与参数估计可参照Rabiner的教程\cite{rabiner1989hmm}。下面将通用过滤递推具体化为有限求和。
设$x_{1:T}$为一段设备观测，$z_t\in\{1,\ldots,K\}$为时刻$t$的隐状态。HMM采用一阶马尔可夫性和观测条件独立性：
\begin{align}
 p(z_t\mid z_{1:t-1})&=p(z_t\mid z_{t-1}),\\
 p(x_t\mid z_{1:T},x_{1:t-1})&=p(x_t\mid z_t).
\end{align}
记初始概率为$\pi_k$，转移概率为$A_{jk}=p(z_t=k\mid z_{t-1}=j)$，观测密度或概率为$b_k(x)=p(x\mid z_t=k)$。其中$\pi\in\R^K$、$A\in\R^{K\times K}$，行索引$j$是出发状态，列索引$k$是到达状态；离散观测的$b_k$为概率，连续观测时为密度。联合分布分解为
\begin{equation}
 p(x_{1:T},z_{1:T})=\pi_{z_1}b_{z_1}(x_1)
 \prod_{t=2}^{T}A_{z_{t-1},z_t}b_{z_t}(x_t).\label{eq:bayes-hmm-joint}
\end{equation}
\begin{figure}[htbp]
\centering
\begin{tikzpicture}[>=Latex,
 state/.style={circle,draw,minimum size=10mm,fill=white},
 obs/.style={circle,draw,minimum size=10mm,fill=blue!12},
 every node/.style={font=\small}]
 \foreach \t in {1,2,3,4}{
  \node[state] (z\t) at (2.4*\t,0) {$z_{\t}$};
  \node[obs] (x\t) at (2.4*\t,-1.6) {$x_{\t}$};
  \draw[->,thick] (z\t) -- (x\t);
 }
 \foreach \a/\b in {1/2,2/3,3/4}{\draw[->,thick] (z\a) -- (z\b);}
 \node[above=3mm of z1] {$\pi$};
 \node at (6,0.65) {状态转移：$p(z_t\mid z_{t-1})$};
 \node at (6,-2.45) {观测生成：$p(x_t\mid z_t)$};
\end{tikzpicture}
\caption{HMM 的有向图结构（以四个时刻为例）。白色圆为隐状态，蓝色圆为已观测变量；横向箭头表示状态转移，向下箭头表示观测生成。箭头表示概率分解关系，不自动代表因果关系。}
\label{fig:bayes-hmm-structure}
\end{figure}
图\ref{fig:bayes-hmm-structure}可以直接读成式\eqref{eq:bayes-hmm-joint}的联合分布：每个节点贡献一个给定父节点的条件分布，首个状态贡献初始分布。给定$z_t$后，$x_t$与其他时刻的变量条件独立；但只给定观测而未固定隐状态时，各时刻观测通常仍然相关。图中的向下箭头是模型定义的生成方向，推断时则可以从观测反过来更新隐状态的概率。

每行转移概率之和为一。多通道观测可以用联合高斯分布或混合分布描述，HMM 并不要求同一时刻的各传感器通道独立。真正受限制的是跨时刻的依赖：一旦给定状态序列，观测之间便被假定条件独立。

以“正常、退化、故障”三个状态为例，转移矩阵可以表达状态持续性。禁止某条转移必须有物理依据，例如只有在明确排除维修过程的运行片段内，才可能限制故障状态不能直接返回正常。无监督学习得到的三个隐状态也不会自动具有这三个语义，需要通过标注、特征分布或专家检查建立对应关系。

\paragraph{过滤、平滑与最优路径。}
在线维护需要$p(z_t\mid x_{1:t})$，称为过滤；事后分析使用$p(z_t\mid x_{1:T})$，称为平滑。后者利用未来观测，不能直接作为实时告警的性能。若需要解释整段运行过程，还可以寻找联合概率最大的状态路径。这三个问题的输出不同，即使使用同一个 HMM，也不能混用结果。

直接枚举所有状态路径需要考虑$K^T$种配置。前向算法利用链结构复用公共计算，令$\alpha_t(k)=p(x_{1:t},z_t=k)$，则
\begin{align}
 \alpha_1(k)&=\pi_kb_k(x_1),\\
 \alpha_t(k)&=b_k(x_t)\sum_{j=1}^{K}\alpha_{t-1}(j)A_{jk},\\
 p(z_t=k\mid x_{1:t})&=\frac{\alpha_t(k)}{\sum_j\alpha_t(j)}.
\end{align}
最后一时刻的前向量求和即为$p(x_{1:T})$。每步对$K^2$个状态对计算，总复杂度为$O(TK^2)$。长序列中概率连乘会下溢，应逐步缩放或在对数域使用 log-sum-exp，并保留缩放因子以恢复序列似然。

后向量$\beta_t(j)=p(x_{t+1:T}\mid z_t=j)$满足
\begin{equation}
 \beta_T(j)=1,\qquad
 \beta_t(j)=\sum_{k=1}^{K}A_{jk}b_k(x_{t+1})\beta_{t+1}(k).
\end{equation}
于是平滑概率为$\gamma_t(k)=\alpha_t(k)\beta_t(k)/p(x_{1:T})$。前向量汇总过去证据，后向量汇总未来证据，这也给出了链式消息传递的具体形式。

Viterbi 算法将前向递推中的求和替换为最大值。在对数域中，
\begin{equation}
 \delta_t(k)=\log b_k(x_t)+\max_j\{\delta_{t-1}(j)+\log A_{jk}\},
 \qquad \delta_1(k)=\log\pi_k+\log b_k(x_1).
\end{equation}
保存每步最大值对应的前驱，便可从终点回溯路径。逐时刻选择边缘概率最大的状态与 Viterbi 路径通常不同，前者优化逐点分类，后者优化整条路径的联合概率。在存在禁止转移时，逐点选择甚至可能产生不合法路径。

\paragraph{两状态、两时刻的过滤手算。}
设状态为正常$N$、退化$D$，$\pi=(0.8,0.2)$，按$N,D$排列的转移矩阵为
\begin{equation}
 A=\begin{bmatrix}0.9&0.1\\0.2&0.8\end{bmatrix},\qquad
 P(\text{高振动}\mid N)=0.1,\quad P(\text{高振动}\mid D)=0.7.
\end{equation}
首次看到高振动时，$\alpha_1=(0.08,0.14)$，所以退化过滤概率为$0.14/0.22=7/11$。第二次仍为高振动，先转移再乘观测概率，得到$\alpha_2(N)=(0.08\times0.9+0.14\times0.2)\times0.1=0.0100$，$\alpha_2(D)=(0.08\times0.1+0.14\times0.8)\times0.7=0.0840$。新的退化概率为$0.084/0.094\approx0.894$。

若事后回看第一个时刻，$\beta_1(N)=0.9\times0.1+0.1\times0.7=0.16$，$\beta_1(D)=0.2\times0.1+0.8\times0.7=0.58$，所以退化平滑概率为$0.14\times0.58/0.094\approx0.864$。它不同于当时可算的$7/11$，因为第二次高振动已经提供了未来证据。此处$\alpha_t,\beta_t,\gamma_t$各有$K$个分量，$A$为$K\times K$矩阵，$\xi_t$也为$K\times K$相邻状态概率表。

\paragraph{前向后向递推的推导与稳定实现。}
固定当前状态$k$，对前驱$j$边缘化：
\begin{align}
 p(x_{1:t},z_t=k)
 &=\sum_jp(x_t\mid z_t=k)p(z_t=k\mid z_{t-1}=j)
             p(x_{1:t-1},z_{t-1}=j)\\
 &=b_k(x_t)\sum_jA_{jk}\alpha_{t-1}(j).
\end{align}
条件独立允许把观测因子提到求和外。相似地，从$z_{t+1}$边缘化得到后向递推；给定$z_t$后过去和未来的观测分离，因此$\alpha_t(k)\beta_t(k)$就是$p(x_{1:T},z_t=k)$。求和表示保留所有可能路径，Viterbi的最大值则丢弃每个终点下非最优前缀；链结构保证最优前缀可扩展为最优整段路径。

数值上令$h_{t-1}(j)=p(z_{t-1}=j\mid x_{1:t-1})$，计算
\begin{equation}
 u_t(k)=b_k(x_t)\sum_jh_{t-1}(j)A_{jk},\quad
 c_t=\sum_ku_t(k),\quad h_t(k)=u_t(k)/c_t.
\end{equation}
第一步用$u_1(k)=\pi_kb_k(x_1)$。于是$c_t=p(x_t\mid x_{1:t-1})$且$\log p(x_{1:T})=\sum_t\log c_t$；缩放后的后向递推对应除以$c_{t+1}$。精确运算下$c_t=0$表示当前证据不在模型支持内；浮点计算中的零也可能来自观测密度或乘积下溢，应先用对数域区分二者，再检查零转移与观测支持。对数域的$\log\sum_j e^{a_j}=m+\log\sum_je^{a_j-m}$（$m=\max_ja_j$有限）提供稳定计算；若所有$a_j=-\infty$，结果直接取$-\infty$，不能计算$-\infty-(-\infty)$。

\paragraph{从隐状态推断到参数学习。}
隐状态未标注时，Baum--Welch 算法以 Expectation Maximization (EM) 交替估计状态的后验占用和模型参数。E 步计算$\gamma_t(j)$以及相邻状态的后验概率
\begin{equation}
 \xi_t(j,k)=\frac{\alpha_t(j)A_{jk}b_k(x_{t+1})\beta_{t+1}(k)}{p(x_{1:T})}.
\end{equation}
对一条序列，M 步的转移更新为
\begin{equation}
 A_{jk}^{\mathrm{new}}=
 \frac{\sum_{t=1}^{T-1}\xi_t(j,k)}{\sum_{t=1}^{T-1}\gamma_t(j)}.
\end{equation}
初始分布更新为$\pi_k^{\mathrm{new}}=\gamma_1(k)$，观测分布则按$\gamma_t(k)$进行加权估计。多条独立运行序列需要分别推断后再汇总计数，不能把不同设备的尾部与头部连接成一次状态转移。EM 依赖初始化，只保证在精确 E、M 步下不降低训练似然，不能保证找到全局最优参数。

\paragraph{EM的期望计数更新。}
取$q(z)=p(z\mid x,\theta^{\rm old})$，由Jensen不等式，
\begin{equation}
 \log p_\theta(x)\geq\E_q\log p_\theta(x,z)-\E_q\log q(z).
\end{equation}
在旧参数处等号成立；M步保持$q$固定并提高右端，便不降低观测似然。转移相关部分是$\sum_{jk}N_{jk}\log A_{jk}$，其中$N_{jk}=\sum_{t<T}\xi_t(j,k)$。对每行加入约束乘子$\lambda_j(\sum_kA_{jk}-1)$，一阶条件$N_{jk}/A_{jk}+\lambda_j=0$与行和为一给出$A_{jk}=N_{jk}/\sum_kN_{jk}$。又因$\sum_k\xi_t(j,k)=\gamma_t(j)$，便得到前述更新。

若$x_t\in\R^d$且$b_k=\mathcal N(\mu_k,\Sigma_k)$，高斯对数密度的加权一阶条件给出
\begin{equation}
 \mu_k^{\rm new}=\frac{\sum_t\gamma_t(k)x_t}{N_k},\qquad
 \Sigma_k^{\rm new}=\frac{\sum_t\gamma_t(k)(x_t-\mu_k^{\rm new})(x_t-\mu_k^{\rm new})^{\mathsf T}}{N_k},
 \quad N_k=\sum_t\gamma_t(k).
\end{equation}
这里$\mu_k\in\R^d$、$\Sigma_k\in\R^{d\times d}$。若状态从未被占用，分母为零，参数不可由本轮数据识别；协方差也可能退化，需预先规定先验、协方差结构或数值下界，而非把无定义更新当作成功训练。约束或近似M步需要实际检查目标是否改善。

若状态转移已经观测到，可以为$A_{j,:}$设置 Dirichlet 先验，其后验参数等于先验参数加转移计数。状态未知时，转移计数本身也不确定。EM 的期望计数可用于点估计，但将其代入 Dirichlet 分布并不自动得到精确的参数后验，完整贝叶斯学习需要联合处理状态和参数的不确定性。

标准齐次HMM的自转移还隐含持续时间假设。若$A_{kk}<1$，进入状态$k$后停留$d$步的概率为
\begin{equation}
 p(D=d\mid k)=A_{kk}^{d-1}(1-A_{kk}),\qquad d=1,2,\ldots.
\end{equation}
这是一种无记忆的几何分布。若退化持续时间明显集中在某个范围，需要重新检查状态定义或显式建立持续时间模型。若将HMM所有转移行设为同一个$\pi$，并令初始分布也为$\pi$，各时刻隐状态就独立同分布。高斯观测下，边缘化各状态便恢复独立样本的GMM。这个特例说明，从混合模型到序列模型所增加的结构是隐状态之间的依赖。

\section{局部条件模型与MEMM}
\subsection{为什么直接建模标签的条件分布}
HMM同时描述状态和观测的联合分布。当目标仅是识别状态，而输入包含频谱、载荷变化和重叠窗口特征时，为所有输入建立合适的生成分布可能很困难。条件模型直接学习$p(z\mid x)$，允许使用相关观测作为输入，无需假定这些输入彼此独立，也无需为$p(x)$建模。这样处理的是观测特征的使用限制，并没有显式建立观测之间的联合分布。

最大熵马尔可夫模型（Maximum Entropy Markov Model，MEMM）保留标签的有向链，以条件最大熵分类器描述每步转移\cite{mccallum2000memm}：
\begin{equation}
 p_\theta(z_t=k\mid z_{t-1}=j,x)=
 \frac{\exp\{\theta^{\mathsf T}f(j,k,x,t)\}}
 {\sum_{k'}\exp\{\theta^{\mathsf T}f(j,k',x,t)\}}.
\end{equation}
其中$f$为特征向量，$x$表示当前任务允许使用的观测范围。最大熵的名称来自特征期望约束下的条件熵最大化，其解具有上述指数族形式。令$z_0$为固定起始状态，序列条件概率为各步条件概率之积。对于完整标注序列，训练可转化为以真实前驱状态为条件的局部分类；预测时前驱状态未知，仍需进行链上推断。

\begin{figure}[htbp]
\centering
\begin{tikzpicture}[>=Latex,
 state/.style={circle,draw,minimum size=10mm},
 every node/.style={font=\small}]
 \node[draw,rounded corners,fill=blue!12,minimum width=9cm,minimum height=9mm] (x) at (5,1.8) {已给定的可用观测 $x$（不为其建立生成分布）};
 \node[state,fill=gray!15] (z0) at (0,0) {$z_0$};
 \foreach \t in {1,2,3,4}{\node[state] (z\t) at (2.4*\t,0) {$z_{\t}$};}
 \foreach \a/\b in {0/1,1/2,2/3,3/4}{\draw[->,thick] (z\a) -- (z\b);}
 \foreach \t in {1,2,3,4}{\draw[->,blue!65!black] (x.south -| z\t.north) -- (z\t.north);}
 \node[align=center] at (4.8,-1.15) {每一步固定前驱 $z_{t-1}$，只对后继 $z_t$ 归一化\\$\displaystyle \sum_k p_\theta(z_t=k\mid z_{t-1},x)=1$};
\end{tikzpicture}
\caption{MEMM 的条件有向结构。灰色节点为固定起始状态，观测同时参与各步局部分类；与 HMM 相比，模型不再使用由状态指向观测的生成边。}
\label{fig:bayes-memm-structure}
\end{figure}
图\ref{fig:bayes-memm-structure}中，每个待预测标签的输入来自前驱标签与已给定观测，因此整条路径概率是这些局部条件概率的乘积。顶部的$x$可以包含任务允许的多个时刻和重叠特征，并不意味着在线任务能够读取未来。横向箭头保留了标签顺序依赖，却没有取消每一步各自归一化的限制；这正是后文标签偏置的结构来源。

\paragraph{最大熵为何导出softmax。}
把前驱与可用输入合称上下文$c$，令经验上下文权重为$\hat p(c)$，特征$\phi(c,k)\in\R^p$。在归一化约束$\sum_kq(k\mid c)=1$及特征期望约束$\sum_{c,k}\hat p(c)q(k\mid c)\phi(c,k)=\bar\phi$下，最大化条件熵$-\sum_{c,k}\hat p(c)q(k\mid c)\log q(k\mid c)$。对$q(k\mid c)$求拉格朗日一阶条件得到
\begin{equation}
 -\log q(k\mid c)-1+\theta^{\mathsf T}\phi(c,k)+\lambda(c)=0.
\end{equation}
指数化后，归一化消去$\lambda(c)$，即得$q(k\mid c)=\exp(\theta^{\mathsf T}\phi(c,k))/\sum_{k'}\exp(\theta^{\mathsf T}\phi(c,k'))$。这要求约束可行；特征完全分离时有限最优权重未必存在，通常以正则化松弛精确矩约束。名称中的熵指这个构造原则，不表示模型输出必须高熵。

MEMM 可以使用相关、重叠的输入特征，不需要描述它们的联合分布，但每个前驱状态都单独归一化，这带来了标签偏置。若某个状态只有一个允许的后继，该转移的条件概率恒为一，无论对应特征得分多低，都不能借助该步降低路径概率。更一般地，后继较少或转移分布较集中的状态可能获得结构性优势。问题来自局部归一化对路径评分的约束，并非简单的类别数量不平衡。

\section{无向图与条件随机场}
\subsection{无向图怎样表达条件独立}
为了让整条标签路径共同参与归一化，可以用无向图构造条件分布。先讨论一般无向图。无向边不区分父节点和子节点，表达变量之间允许的直接相互作用。给定其他所有变量后，没有直接边的两点在满足图的马尔可夫性质时条件独立。

无向图的分离规则比d-分离简单。若从节点集合$A$到$B$的每条路径都经过条件集合$S$，则$S$分离$A,B$，图的全局马尔可夫性质给出
\begin{equation}
 X_A\perp X_B\mid X_S.
 \label{eq:bayes-undirected-separation}
\end{equation}
这里$A,B,S$两两不相交。实际操作可以删除$S$及其关联边，检查剩余图中$A$与$B$是否仍连通。无向图没有碰撞节点，因此没有“观察碰撞节点反而打开路径”的判断规则。链$A-B-C$中，给定$B$就分离两端。如果另有路径$A-D-C$，还需要阻断这条路径。

\subsection{从团与势函数构造联合分布}
马尔可夫随机场（Markov Random Field，MRF）用无向图表达上述独立关系。团（Clique）是节点两两相连的集合，极大团是不能再加入其他节点而仍保持两两相连的团。对有限离散变量，严格正的联合分布满足图的全局马尔可夫性质，当且仅当它能按团因子分解。这是Hammersley--Clifford定理的常用形式。

选取覆盖所需相互作用的团集合$\mathcal C$，分布写成
\begin{equation}
 p(x)=\frac{1}{Z}\prod_{C\in\mathcal C}\psi_C(x_C),\qquad
 Z=\sum_x\prod_{C\in\mathcal C}\psi_C(x_C).
 \label{eq:bayes-undirected-factorization}
\end{equation}
$x_C$表示团中变量的取值，势函数（Potential Function）$\psi_C$表示这些取值的相容程度，配分函数（Partition Function）$Z$对全部配置归一化。势函数无需单独归一化，也不是某个方向的条件概率。连续变量把对应求和改成积分，并要求$0<Z<\infty$。连续版本的分解定理还需要相应的密度与正则条件。

链$A-B-C$的极大团为$\{A,B\},\{B,C\}$，可以写成$p(a,b,c)=\psi_{AB}(a,b)\psi_{BC}(b,c)/Z$。固定$b$后，分子分为关于$a$与$c$的两项，分母也分解成两个和的乘积，因此$p(a,c\mid b)=p(a\mid b)p(c\mid b)$。这说明团分解如何产生条件独立。

三角形$A-B-C-A$则包含三元极大团。一般需要$\psi_{ABC}(a,b,c)$，只乘三个二元势会额外限制分布，不能表达所有三元相互作用。若存在零概率配置，按非负势定义分布仍然可行，但不能直接使用依赖严格正性的上述等价定理。以极大团分解时，较小团的势可以吸收到包含它的极大团中，因子表示也不唯一。

\begin{figure}[htbp]
\centering
\begin{tikzpicture}[every node/.style={font=\small},state/.style={circle,draw,minimum size=9mm}]
 \node[state] (a) at (0,0) {$A$};
 \node[state,fill=orange!18] (b) at (1.8,0) {$B$};
 \node[state] (c) at (3.6,0) {$C$};
 \draw[thick] (a)--node[above] {$\psi_{AB}$}(b);
 \draw[thick] (b)--node[above] {$\psi_{BC}$}(c);
 \node[align=center] at (1.8,-1.4) {链：两个二元极大团\\$p(a,b,c)=\psi_{AB}(a,b)\psi_{BC}(b,c)/Z$\\给定$B$分离$A,C$};
 \node[state] (ta) at (6.5,0) {$A$};
 \node[state] (tb) at (8.1,1.2) {$B$};
 \node[state] (tc) at (9.7,0) {$C$};
 \draw[thick] (ta)--(tb)--(tc)--(ta);
 \node[align=center] at (8.1,-1.4) {三角形：一个三元极大团\\$p(a,b,c)=\psi_{ABC}(a,b,c)/Z$\\给定$B$仍有$A-C$直接边};
\end{tikzpicture}
\caption{无向图的分离与团分解。链中的橙色节点表示所给条件；三角形的极大团包含全部三个节点，不能仅凭有三条边就假定只有二元相互作用。两种模型各有自己的配分函数$Z$。}
\label{fig:bayes-clique-factorization}
\end{figure}

图\ref{fig:bayes-clique-factorization}对照了前面的链与三角形。先找极大团，再给每个团定义势函数，最后对所有变量配置归一化。节点是否被条件化与节点是否属于隐变量，是两个不同问题。

\subsection{从局部归一化改为全局归一化}
条件随机场（Conditional Random Field，CRF）对每个给定的观测$x$，用无向图描述标签$Z$的条件分布\cite{lafferty2001crf}：
\begin{equation}
 p_\theta(z\mid x)=\frac{1}{Z_\theta(x)}
 \prod_{C\in\mathcal C}\psi_C(z_C;x),\qquad
 Z_\theta(x)=\sum_z\prod_{C\in\mathcal C}\psi_C(z_C;x).
\end{equation}
$x$作为势函数的已知条件，可以包含多个相关观测。CRF无需对$x$建立生成分布。这里的无向图描述的是给定$x$后的标签关系。

若标签节点构成一条链，就得到线性链条件随机场（Linear-Chain CRF）。标签在训练时通常已观测，预测时未知，因此“线性链”由输出变量的图结构决定，不要求训练标签是隐变量。若训练时也有未观测节点，需要另行边缘化，得到带隐变量的条件模型。

线性链CRF仍保留无向链的局部马尔可夫性质：给定$x$和相邻标签后，一个内部标签与其余标签条件独立。它允许势函数依赖位置和整段可用观测，不必使用齐次的生成转移概率。这不等于取消了所有马尔可夫关系。齐次性、图的条件独立性和归一化方式是三个不同问题。

线性链CRF把归一化放到整条路径上。定义路径得分
\begin{align}
 F_\theta(z,x)&=\sum_{t=1}^{T}\theta^{\mathsf T}f(z_{t-1},z_t,x,t),\\
 p_\theta(z\mid x)&=\frac{\exp F_\theta(z,x)}{Z_\theta(x)},\\
 Z_\theta(x)&=\sum_{z'_{1:T}}\exp F_\theta(z',x).
\end{align}
\begin{figure}[htbp]
\centering
\begin{tikzpicture}[
 state/.style={circle,draw,minimum size=9mm},
 factor/.style={rectangle,draw,fill=orange!25,minimum size=7mm,inner sep=1pt},
 every node/.style={font=\small}]
 \node[draw,rounded corners,fill=blue!12,minimum width=10cm,minimum height=9mm] (x) at (5,1.8) {给定观测 $x$：各局部势函数共享的条件};
 \node[state,fill=gray!15] (z0) at (0,0) {$z_0$};
 \foreach \t in {1,2,3,4}{
  \node[state] (z\t) at (2.5*\t,0) {$z_{\t}$};
  \node[factor] (f\t) at (2.5*\t-1.25,0) {$\psi_{\t}$};
  \draw[thick] (f\t) -- (z\t);
  \draw[dashed,blue!65!black] (x.south -| f\t.north) -- (f\t.north);
 }
 \foreach \a/\b in {0/1,1/2,2/3,3/4}{\draw[thick] (z\a) -- (f\b);}
 \node[align=center] at (5,-1.2) {$\psi_t=\exp\{\theta^{\mathsf T}f(z_{t-1},z_t,x,t)\}$\\局部势相乘，再用同一个 $Z_\theta(x)$ 对整条路径归一化};
\end{tikzpicture}
\caption{线性链 CRF 的因子图。圆形表示标签，方形表示相邻标签的势函数；实线无方向，表示变量参与因子。蓝色虚线仅表示势函数依赖已给定观测，不是标签图中的随机变量边或生成箭头。}
\label{fig:bayes-crf-structure}
\end{figure}
图\ref{fig:bayes-crf-structure}把条件分布写成可逐项检查的乘积：每个方形读取相邻两个标签及观测，输出一个非负相容度。固定$z_0$后，第一个因子实际只涉及一个待预测标签。去掉因子节点、将共享因子的标签直接相连，就得到标签的无向链。与 MEMM 不同，$\psi_t$不要求对后继标签求和为一，所有因子共同决定路径得分，最后才由一个全局配分函数归一化。

配分函数$Z_\theta(x)$比较所有允许路径的得分。一个只有唯一后继的状态仍然可以获得较低的路径得分，因此消除了 MEMM 局部归一化造成的上述限制。全局归一化不意味着自动解决标注噪声、分布偏移或特征不足。

对一条标注序列$(x,z)$，带二次正则的负对数条件似然为
\begin{equation}
 \mathcal L(\theta)=-F_\theta(z,x)+\log Z_\theta(x)
 +\frac{\lambda}{2}\|\theta\|_2^2.
\end{equation}
其梯度等于模型期望的特征计数减去真实路径的特征计数，再加$\lambda\theta$。这个形式解释了训练为何需要边缘概率：只找出当前最优路径不足以计算配分函数的梯度。在线性链、有限状态和相邻标签因子下，前向后向算法可在$O(TK^2)$时间内计算配分函数和边缘概率，Viterbi 则计算最优路径。

\paragraph{两条路径怎样体现全局归一化。}
假设只有两条合法路径，得分分别为0和$\log3$，则未归一化权重为1和3，$Z=4$，路径概率为$1/4$和$3/4$。若观测证据使第二条路径增加$\log2$，它的权重变成6，概率成为$6/7$。即使该路径某一步只有一个合法后继，该步得分仍能改变整条路径的权重；MEMM在那个唯一后继处的局部概率却恒为一。这个小例子只解释归一化的区别，并不是说CRF必然取得更高测试准确率。

特征向量$f(j,k,x,t)\in\R^p$可以把“从正常到退化”“当前振动超过阈值”等指示量列成$p$项，$\theta\in\R^p$为其权重，内积得到一个标量得分。指数化把任何实得分变成正权重，最后归一化才得到概率。后文神经编码器只是在学习怎样产生这些局部得分，不改变概率必须归一化这一要求。

\paragraph{配分函数、边缘概率与梯度的完整联系。}
令$\psi_t(j,k)=\exp\{\theta^{\mathsf T}f(j,k,x,t)\}$。固定起点$z_0$，设置$a_0(z_0)=1$、其余为零，递推$a_t(k)=\sum_ja_{t-1}(j)\psi_t(j,k)$，则$Z=\sum_ka_T(k)$。后向量$b_T(k)=1$、$b_{t-1}(j)=\sum_k\psi_t(j,k)b_t(k)$给出边缘
\begin{equation}
 p_\theta(z_{t-1}=j,z_t=k\mid x)
 =\frac{a_{t-1}(j)\psi_t(j,k)b_t(k)}{Z}.
\end{equation}
这些$a,b$是未归一化消息，不是HMM的观测似然参数。定义总特征
\begin{equation}
 \Phi(z,x)=\sum_tf(z_{t-1},z_t,x,t)\in\R^p.
\end{equation}
直接对有限和求导得
\begin{align}
 \nabla_\theta\log Z
 &=\frac{\sum_z e^{\theta^{\mathsf T}\Phi(z,x)}\Phi(z,x)}{Z}
 =\E_{p_\theta(z\mid x)}\Phi(z,x),\\
 \nabla_\theta\mathcal L
 &=-\Phi(z^{\rm obs},x)+\sum_{t,j,k}p_\theta(j,k\mid x)f(j,k,x,t)+\lambda\theta.
\end{align}
再求导得到$\nabla^2\mathcal L=\operatorname{Cov}_{p_\theta}(\Phi)+\lambda I_p$。因此固定特征的线性CRF目标凸，$\lambda>0$时严格凸；若特征由待训练神经网络产生，这一参数凸性便不再成立。禁止路径可赋零势，但至少必须存在一条合法路径，否则$Z=0$使条件分布无定义。

神经网络可以提供 CRF 的局部观测得分。例如，令$s_t(k;x)$为编码器输出，$B_{jk}$为可学习转移得分，则
\begin{equation}
 F_\theta(z,x)=\sum_{t=1}^{T}s_t(z_t;x)+\sum_{t=2}^{T}B_{z_{t-1},z_t}.
\end{equation}
与 HMM 的$A_{jk}$不同，$B_{jk}$是未归一化得分，不要求非负或行和为一。Bidirectional Long Short-Term Memory (BiLSTM)--CRF 和 Transformer--CRF 将特征提取交给神经网络，将标签关系交给结构化输出层。双向编码器适合完整记录的离线标注；在线预警必须限制输入和推断范围，或者明确等待未来观测所引入的延迟。即使局部特征只读取过去，对整段序列进行平滑仍可能利用后续证据。

\subsection{网格中的空间相互作用}
时间链只是一种邻接关系。在缺陷图像中，每个标签可能有多个邻居。Geman与Geman将Gibbs分布和无向图建模用于图像恢复\cite{geman1984mrf}。沿用式\eqref{eq:bayes-undirected-factorization}的团分解，令$E(z)=-\sum_C\log\psi_C(z_C)$，便得到能量形式：
\begin{equation}
 p(z)=\frac{1}{Z}\prod_{C\in\mathcal C}\psi_C(z_C)
 =\frac{1}{Z}\exp\{-E(z)\}.
\end{equation}
\begin{figure}[htbp]
\centering
\begin{tikzpicture}[
 site/.style={circle,draw,minimum size=11mm,fill=white},
 every node/.style={font=\small}]
 \foreach \r in {1,2,3}{
  \foreach \c in {1,2,3}{\node[site] (v\r\c) at (1.8*\c,-1.5*\r) {$z_{\r\c}$};}
 }
 \foreach \r in {1,2,3}{\foreach \a/\b in {1/2,2/3}{\draw[thick] (v\r\a) -- (v\r\b);}}
 \foreach \c in {1,2,3}{\foreach \a/\b in {1/2,2/3}{\draw[thick] (v\a\c) -- (v\b\c);}}
 \node[site,fill=orange!30] at (v22) {$z_{22}$};
 \foreach \n in {12,21,23,32}{\node[site,fill=green!15] at (v\n) {$z_{\n}$};}
 \node[align=left,text width=5.2cm,anchor=west] at (6.3,-3) {橙色：当前标签 $z_{22}$\\[3pt]绿色：四邻域 $N(22)$\\[3pt]白色：其余标签\\[6pt]给定绿色邻居后，中心标签与其余标签条件独立。};
\end{tikzpicture}
\caption{四邻域网格上的成对 MRF。每个节点对应一个像素或区域标签，每条无向边对应一个二元势；不存在箭头，并不表示各标签独立，而是表示相互作用不采用有向生成分解。}
\label{fig:bayes-mrf-structure}
\end{figure}
图\ref{fig:bayes-mrf-structure}展示的是一种具体的网格结构。此处$z_{rc}$的两个下标分别表示行和列。局部马尔可夫性质可写为
\begin{equation}
 p(z_{22}\mid z_{\setminus 22})
 =p(z_{22}\mid z_{12},z_{21},z_{23},z_{32}),
\end{equation}
其中$z_{\setminus 22}$表示除中心外的全部标签。对于图示的一元和二元势分解，固定其他标签后，与$z_{22}$无关的因子在条件概率归一化时相消，因而
\begin{equation}
 p(z_{22}=k\mid z_{\setminus 22})
 =\frac{\psi_{22}(k)\prod_{j\in N(22)}\psi_{22,j}(k,z_j)}
 {\sum_{k'}\psi_{22}(k')\prod_{j\in N(22)}\psi_{22,j}(k',z_j)}.
\end{equation}
这也解释了Gibbs采样为何只需读取当前节点的邻居。不过，局部更新只访问邻域，并不意味着联合分布能够分解成相互独立的节点分布。若这些势还依赖给定图像$x$，图中的标签邻接关系保持不变，就得到前面定义的CRF在网格上的形式。

团$C$内的节点两两相连，势函数$\psi_C$表达局部相容程度，并不是归一化概率。能量$E(z)=-\sum_C\log\psi_C(z_C)$越低，配置概率越高。配分函数$Z$对全部配置求和，保证总概率为一。若使用零势函数表达硬约束，则不能直接套用依赖严格正性的等价结论。

在给定图像$x$的二元缺陷标注中，可以设置
\begin{equation}
 E(z;x)=\sum_i U_i(z_i;x)+\lambda\sum_{(i,j)\in\mathcal E}
 w_{ij}(x)\mathbb 1\{z_i\ne z_j\},\qquad \lambda,w_{ij}(x)\geq0.
\end{equation}
一元项$U_i$衡量局部像素证据，二元 Potts 项惩罚邻居标签不一致。相邻区域越相似，$w_{ij}$可以越大；跨越明显边界时降低权重，可减轻过度平滑。若直接将分类概率转换成一元负对数得分，这构成条件模型的一种参数化，并不自动成为生成似然。

此时$p(z\mid x)\propto\exp\{-E(z;x)\}$也是一个CRF。MRF强调无向马尔可夫依赖，CRF强调给定观测后的结构化条件分布。每个固定$x$下的CRF都按对应无向图分解。线性链CRF是其中计算特别方便的一类，二维网格CRF通常不能直接沿用链式动态规划。

能量中的平滑项可以移除孤立误报，也可能抹掉真实细裂纹，因此应依据缺陷宽度和边界误差选择强度。一般含环图的精确推断可能十分昂贵，变量消元的复杂度随树宽增长。置信传播在树上精确，在含环图上通常只是近似且不保证收敛。均值场以较简单的分布逼近联合依赖，Gibbs 采样则逐变量更新。对于二元、次模的成对能量，图割可求全局最小能量，但这一结论不能推广到任意多类别势函数或边缘概率计算。

\begin{table}[htbp]
\centering
\caption{链式与无向图模型的建模对象和计算边界。}
\label{tab:bayes-graph-models}
\small
\begin{tabular}{m{1.3cm}m{3.1cm}m{3.3cm}m{4.2cm}}
\hline
模型 & 建模对象 & 归一化方式 & 关键限制 \\
\hline
HMM & 联合分布$p(x,z)$ & 转移与观测局部归一化 & 观测假设和状态持续时间需检验 \\
MEMM & 条件分布$p(z\mid x)$ & 每个前驱状态局部归一化 & 标签偏置影响路径评分 \\
链式 CRF & 条件分布$p(z\mid x)$ & 所有标签路径全局归一化 & 训练需计算模型期望 \\
MRF & 无向图上的联合分布 & 全部配置全局归一化 & 一般图中配分函数和推断困难 \\
\hline
\end{tabular}
\end{table}

表\ref{tab:bayes-graph-models}按建模对象、归一化方式和推断边界比较图模型；应先决定需要联合生成还是条件预测，再判断链式动态规划是否适用。
\paragraph{MRF的局部更新与均值场近似。}
对固定其他标签$z_{-i}$，不包含$i$的能量项在条件概率的分子分母中相消，因此
\begin{equation}
 p(z_i=k\mid z_{-i},x)\propto
 \exp\left\{-U_i(k;x)-\lambda\sum_{j:(i,j)\in\mathcal E}w_{ij}(x)\mathbb 1\{k\ne z_j\}\right\}.
\end{equation}
逐点从此分布抽样就是Gibbs更新；若改为取最小局部能量，则变为坐标下降，只能保证能量不增而非全局最优。有限状态、适当遍历性和非周期性下Gibbs链以目标分布为平稳分布，强耦合时仍可能混合极慢。

若用乘积分布$q(z)=\prod_iq_i(z_i)$近似，最小化$\KL(q\|p)$等价于最小化以下目标：
\begin{equation}
 \E_qE+\sum_i\sum_kq_i(k)\log q_i(k).
\end{equation}
固定其他因子并对$q_i$加归一化乘子求导，可得
\begin{equation}
 q_i(k)\propto\exp\{-\E_{q_{-i}}E(k,z_{-i};x)\}
 \propto\exp\{-U_i(k;x)-\lambda\sum_jw_{ij}(1-q_j(k))\}.
\end{equation}
这以邻居概率代替确定标签，但忽略联合相关性；逐坐标精确更新改善变分目标，却不保证找到全局最佳近似。二元成对能量可用图割的关键条件是$V_{ij}(0,0)+V_{ij}(1,1)\leq V_{ij}(0,1)+V_{ij}(1,0)$，非负Potts项满足它。图割解决最大概率配置，不提供配分函数或后验区间。

\paragraph{两个像素的能量可以全部列完。}
设两个相邻标签$z_1,z_2\in\{0,1\}$没有一元偏好，平滑惩罚为$\lambda=\log3$、边权为一。同标签$(0,0),(1,1)$的能量为零、权重为1，不同标签$(0,1),(1,0)$的能量为$\log3$、权重为$1/3$。配分函数为$8/3$，因而$P(0,0)=P(1,1)=3/8$、$P(0,1)=P(1,0)=1/8$。未观察邻居时，像素1的边缘概率为$P(z_1=0)=3/8+1/8=1/2$；已知$z_2=0$时，条件概率为$P(z_1=0\mid z_2=0)=(3/8)/(3/8+1/8)=3/4$。两种配置并列最优，并不使单个像素的边缘判断自动偏向零。在邻居为零时，Gibbs更新以该概率抽零，以$1/4$抽一；若总选择当前条件概率最高的标签，则变成逐坐标寻找局部最优配置，已经不是从后验抽样。配置概率最高与某像素的边缘概率，是两个不同问题。

\paragraph{把图结构带回维护任务。}
假设要从一次运行记录中识别正常、退化与故障阶段，可以先用独立分类器建立逐点基线。若预测频繁跳变，HMM 用转移持续性与观测模型共同解释状态；若观测特征复杂而状态标注充分，MEMM 和 CRF 可直接学习条件边界。若任务转为表面缺陷的空间连通区域，则链结构不再适用，需要结合图像邻接关系定义条件随机场。模型变化由输出依赖决定，而非由名称的新旧决定。

这些比较必须使用相同的可用观测、标签定义和设备划分。逐点准确率之外，还应检查状态切换延迟、阶段持续时间误差和关键故障段漏检。离线平滑结果只能与离线结果比较，在线过滤应另报实际延迟。若由隐状态概率推导未来故障风险，还需要定义故障事件与状态的关系，并验证预测时间窗，不能把“当前处于退化状态”的概率直接解释成“未来二十四小时故障”的概率。

\section{玻尔兹曼机与受限玻尔兹曼机}
\subsection{能量形式与指数族参数化}
前面的无向模型用势函数表达相互作用。对正势函数，取负对数可以把乘积改为能量之和：
\begin{equation}
 E_\theta(x)=-\sum_C\log\psi_C(x_C),\qquad
 p_\theta(x)=\frac{\exp\{-E_\theta(x)\}}{Z_\theta}.
 \label{eq:bayes-gibbs-energy}
\end{equation}
能量是概率权重的参数化，低能量配置具有较高概率。本章把温度常数吸收到参数中，能量因而作为无量纲得分使用。

如果能量关于参数线性，就得到指数族（Exponential Family）的一种常用形式：
\begin{equation}
 p_\theta(x)=h(x)\exp\{\theta^{\mathsf T}T(x)-A(\theta)\},
 \qquad A(\theta)=\log\int h(x)e^{\theta^{\mathsf T}T(x)}\,dx.
 \label{eq:bayes-exponential-family}
\end{equation}
$T(x)$为充分统计量，$\theta$为自然参数，$h(x)$为与参数无关的基准项，$A(\theta)$为对数配分函数。离散变量使用求和。参数必须使归一化量有限。高斯分布、伯努利分布和固定特征的对数线性模型都可写成指数族形式。边的方向决定依赖如何组织，指数族则规定分布怎样参数化，二者不是同一类结构限制。任意参数化的能量模型也不一定是关于所用参数的有限维指数族。

\subsection{玻尔兹曼机怎样引入隐变量}
若只有观测节点，表示复杂的观测联合分布可能需要较多相互作用项。可以在无向图中加入隐节点，让多个观测共享未观测因素。玻尔兹曼机（Boltzmann Machine，BM）是一类具有二元节点和成对相互作用的无向能量模型\cite{bishop2006pattern}。它可以含隐节点，但在任意无向图中加入隐节点并不会自动得到玻尔兹曼机，还需要规定节点取值和能量形式。

\begin{figure}[htbp]
\centering
\begin{tikzpicture}[every node/.style={font=\small},state/.style={circle,draw,minimum size=10mm}]
 \node[state,fill=blue!12] (v1) at (0,0) {$v_1$};
 \node[state,fill=blue!12] (v2) at (2,1.4) {$v_2$};
 \node[state,fill=blue!12] (v3) at (2,-1.4) {$v_3$};
 \node[state,fill=orange!18] (h1) at (4,0.9) {$h_1$};
 \node[state,fill=orange!18] (h2) at (4,-0.9) {$h_2$};
 \draw (v1)--(v2)--(v3)--(v1);
 \draw (h1)--(h2);
 \draw (v1)--(h1); \draw (v1)--(h2); \draw (v2)--(h1);
 \draw (v3)--(h2); \draw (v2)--(h2);
 \node[align=left,anchor=west] at (5.1,0) {一般BM：\\可见节点与隐节点之间\\以及同层节点之间都允许连边};
\end{tikzpicture}
\caption{一般玻尔兹曼机的无向图。节点为二元随机变量，边表示成对相互作用。隐节点之间及同一层可见节点之间的连边，使条件分布通常不再逐节点独立。}
\label{fig:bayes-bm-structure}
\end{figure}

图\ref{fig:bayes-bm-structure}中的边没有方向。它们共同决定能量函数中的相互作用项。RBM通过删除所有同层边得到后面的二部结构。

将可见变量$v\in\{0,1\}^d$与隐变量$h\in\{0,1\}^m$拼成$s=(v^{\mathsf T},h^{\mathsf T})^{\mathsf T}$。可见表示训练数据中可观测，隐变量表示未直接观测的内部节点。一般二元玻尔兹曼机的能量为
\begin{equation}
 E(s)=-\sum_a c_as_a-\sum_{a<b}J_{ab}s_as_b,
 \qquad p(s)=\frac{e^{-E(s)}}{\sum_{s'}e^{-E(s')}}.
\end{equation}
$c_a$是节点偏置，$J_{ab}$是对称相互作用参数，每对节点只计一次。$J_{ab}=0$对应缺少该相互作用边。非零参数可以为正或负，分别改变共同取一的相对权重。有限二元配置和有限参数保证配分函数正且有限。

完整变量$s$的分布属于指数族，充分统计量为$s_a$与$s_as_b$。若只观察$v$，需要求和$p(v)=\sum_h p(v,h)$。这个边缘分布仍是归一化分布，但通常不再保留关于原参数的简单线性充分统计量形式。一般BM允许可见层内部、隐层内部及两层之间连边，隐节点之间的相互作用使$p(h\mid v)$的计算也可能困难。

\subsection{受限玻尔兹曼机的二部图与能量}
受限玻尔兹曼机（Restricted Boltzmann Machine，RBM）删除同一层内的全部连边，只保留可见层与隐层之间的边。这个限制形成二部无向图。它改变了条件独立关系，使给定一层时另一层的节点可以分别计算。

\begin{figure}[htbp]
\centering
\begin{tikzpicture}[every node/.style={font=\small},
 state/.style={circle,draw,minimum size=10mm}]
 \foreach \j in {1,2,3}{\node[state] (h\j) at (2*\j,1.6) {$h_{\j}$};}
 \foreach \i in {1,2,3,4}{\node[state,fill=blue!12] (v\i) at (1.6*\i,0) {$v_{\i}$};}
 \foreach \i in {1,2,3,4}{\foreach \j in {1,2,3}{\draw (v\i)--(h\j);}}
 \node[anchor=west] at (7,1.6) {隐层};
 \node[anchor=west] at (7,0) {可见层};
\end{tikzpicture}
\caption{RBM的二部无向图。图示为全连接的两层结构，也可以删去部分跨层边。每层内部无边，给定另一层后本层节点条件独立。}
\label{fig:bayes-rbm-structure}
\end{figure}

二元RBM沿用$v\in\{0,1\}^d,h\in\{0,1\}^m$，能量与联合分布为
\begin{align}
 E_\theta(v,h)&=-a^{\mathsf T}v-b^{\mathsf T}h-v^{\mathsf T}Wh,\\
 p_\theta(v,h)&=\frac{\exp\{a^{\mathsf T}v+b^{\mathsf T}h+v^{\mathsf T}Wh\}}{Z_\theta},\\
 Z_\theta&=\sum_{v\in\{0,1\}^d}\sum_{h\in\{0,1\}^m}e^{-E_\theta(v,h)}.
 \label{eq:bayes-rbm-joint}
\end{align}
$a\in\R^d,b\in\R^m$分别为可见和隐节点偏置，$W\in\R^{d\times m}$为跨层权重。联合分布的充分统计量是$v_i,h_j,v_ih_j$。图\ref{fig:bayes-rbm-structure}对应的团至多包含一对跨层节点，因此能量中没有$v_iv_{i'}$或$h_jh_{j'}$项。

\subsection{为什么两层的条件分布容易计算}
固定$v$后，把所有包含$h$的项按隐节点整理：
\begin{equation}
 p(h\mid v)\propto\exp\left\{\sum_{j=1}^m h_j
 \left(b_j+\sum_iW_{ij}v_i\right)\right\}
 =\prod_j\exp\{h_j u_j(v)\},
\end{equation}
其中$u_j(v)=b_j+\sum_iW_{ij}v_i$。每个$h_j$只有$0,1$两个取值，归一化量为$1+e^{u_j(v)}$，所以
\begin{align}
 p(h\mid v)&=\prod_jp(h_j\mid v),&
 p(h_j=1\mid v)&=\sigma\left(b_j+\sum_iW_{ij}v_i\right),\\
 p(v\mid h)&=\prod_ip(v_i\mid h),&
 p(v_i=1\mid h)&=\sigma\left(a_i+\sum_jW_{ij}h_j\right),
 \label{eq:bayes-rbm-conditionals}
\end{align}
$\sigma(u)=1/(1+e^{-u})$为逻辑函数。条件化以后，与待更新层无关的项和全局配分函数都被约掉。也可以直接从图的分离看出：给定全部可见节点就阻断了隐节点之间的所有路径，反向同理。

这种独立性依赖于给定另一层。隐节点未给定时，可见节点仍可能通过共同隐节点相关，因此RBM可以描述相关数据。特征表示常取$\E[h\mid v]$，其第$j$项就是上式的激活概率。它是后验期望，不等于已经观测到的真实物理因素，也不保证每个隐节点都有单一语义。

\subsection{可见分布与一个可枚举的例子}
由于隐节点在给定$v$后独立，可以精确求和消去隐层：
\begin{align}
 \sum_h e^{-E(v,h)}
 &=e^{a^{\mathsf T}v}\prod_j\left(1+e^{b_j+\sum_iW_{ij}v_i}\right),\\
 p(v)&=\frac{e^{-\mathcal F(v)}}{Z_\theta},\qquad
 \mathcal F(v)=-a^{\mathsf T}v-\sum_j\log\left(1+e^{b_j+\sum_iW_{ij}v_i}\right).
 \label{eq:bayes-rbm-free-energy}
\end{align}
$\mathcal F(v)$称为自由能，已经对隐变量求和。它与固定某个$h$时的联合能量不同。对数中的求和来源于隐节点的两个可能取值。隐层求和容易，并不意味着$Z_\theta=\sum_v e^{-\mathcal F(v)}$容易，因为可见配置仍有$2^d$个。

设两个可见节点、一个隐节点，全部偏置为零，两个跨层权重均为$\log2$。则
\begin{equation}
 \sum_{h=0}^1e^{-E(v,h)}=1+2^{v_1+v_2}.
\end{equation}
$(v_1,v_2)=(0,0),(1,0),(0,1),(1,1)$的未归一化权重分别为$2,3,3,5$，配分函数为$13$。于是$p(v_1=1,v_2=1)=5/13$，而$p(v_1=1)p(v_2=1)=(8/13)^2$，二者不等。图中可见层没有直接边，边缘上仍然相关。给定$h=0$时，两可见节点独立且取一的概率都是$1/2$。给定$h=1$时，它们仍独立，取一的概率都是$2/3$。

\subsection{最大似然为什么需要两类期望}
训练数据只提供$v$。对单个样本求对数似然梯度，利用$\partial\log Z_\theta/\partial W_{ij}=\E_{p_\theta(v,h)}[v_ih_j]$，得到
\begin{equation}
 \frac{\partial\log p_\theta(v)}{\partial W_{ij}}
 =v_i\,p_\theta(h_j=1\mid v)
 -\E_{p_\theta(v',h')}[v_i'h_j'].
 \label{eq:bayes-rbm-gradient}
\end{equation}
第一项固定真实样本，对隐节点后验求平均，称为正相。第二项对当前模型的联合分布求平均，称为负相。梯度提高数据支持的相互作用，同时减去模型已经产生的相互作用。偏置梯度类似，分别比较可见或隐节点的数据期望与模型期望。对多个样本训练时，再对正相作样本平均。

正相可由式\eqref{eq:bayes-rbm-conditionals}精确计算。负相通常需要采样：从给定$v$抽取整层$h$，再从给定$h$抽取整层$v$，交替进行块Gibbs采样。在有限二元状态和有限参数下，各条件概率严格位于零与一之间，这条链具有目标联合分布作为平稳分布。但有限步样本未必接近平稳分布，强耦合时混合可能很慢。

对比散度（Contrastive Divergence，CD）从真实样本$v^{(0)}$开始，执行$k$轮交替采样，用所得$v^{(k)}$近似负相，称为CD-$k$。常用的权重更新近似为
\begin{equation}
 \Delta W_{ij}=\eta\left[
 v_i^{(0)}p(h_j=1\mid v^{(0)})
 -v_i^{(k)}p(h_j=1\mid v^{(k)})\right],
\end{equation}
其中$\eta>0$为学习率，批训练使用这些量的平均。用条件期望替代末端隐节点抽样可减少部分采样噪声。短链通常没有达到平稳分布，所以CD-$k$是有偏的最大似然梯度近似，不保证每一步提高真实似然。持续对比散度（Persistent Contrastive Divergence，PCD）保留链的状态供下一次更新继续使用，也仍需检查混合质量。

\subsection{连续观测与适用边界}
上述伯努利可见节点适合二元特征，不能直接把实数温度代入后仍解释为伯努利观测。固定尺度的高斯--伯努利RBM可采用
\begin{equation}
 E(v,h)=\sum_i\frac{(v_i-a_i)^2}{2s_i^2}
 -\sum_jb_jh_j-\sum_{i,j}\frac{v_i}{s_i}W_{ij}h_j,
 \qquad s_i>0.
\end{equation}
这里$v\in\R^d$，$h\in\{0,1\}^m$，$s_i$是可见节点的尺度。配方后可得
\begin{align}
 p(v_i\mid h)&=\mathcal N\left(a_i+s_i\sum_jW_{ij}h_j,s_i^2\right),\\
 p(h_j=1\mid v)&=\sigma\left(b_j+\sum_iW_{ij}v_i/s_i\right).
\end{align}
\begin{figure}[htbp]
\centering
\begin{tikzpicture}[every node/.style={font=\small},hidden/.style={circle,draw,minimum size=10mm},visible/.style={ellipse,draw,fill=blue!12,minimum width=22mm,minimum height=10mm}]
 \node[hidden] (h1) at (1,1.6) {$h_1$};
 \node[hidden] (h2) at (4,1.6) {$h_2$};
 \node[visible] (v1) at (0,0) {$v_1\in\R$};
 \node[visible] (v2) at (2.5,0) {$v_2\in\R$};
 \node[visible] (v3) at (5,0) {$v_3\in\R$};
 \foreach \i in {1,2,3}{\foreach \j in {1,2}{\draw (v\i)--(h\j);}}
 \node[anchor=west,align=left] at (6.3,0.8) {隐节点：伯努利条件分布\\可见节点：高斯条件分布\\两层内部仍然无边};
\end{tikzpicture}
\caption{高斯--伯努利RBM。椭圆强调连续可见节点，圆形表示二元隐节点。节点取值改变，二部图的条件独立结构保持不变。}
\label{fig:bayes-gaussian-rbm}
\end{figure}

图\ref{fig:bayes-gaussian-rbm}对照二元RBM显示了连续可见节点的变化。这个变体保留二部图及两层条件独立性，改变的是可见节点的取值与条件分布。二次项保证固定正尺度下连续积分有限。连续特征需处理尺度差异，也需验证高斯条件观测假设。

RBM可以用于学习特征、补全缺失变量和生成样本。缺失值补全时固定已观测节点，更新缺失节点和隐节点，不能把固定观测也重新生成后当作原证据。重构误差小并不保证似然高、样本覆盖充分或异常检测有效。RBM的主要计算收益是两层条件分布容易求解，联合归一化与长期采样仍可能困难。

\section{概率图模型的比较维度}
前面的模型通过改变具体假设得到不同分解。比较时应逐项问清以下问题。参数学习和隐变量推断还要与图结构分开考虑。前者估计分布参数，后者在参数给定时计算未知变量的分布。
\begin{enumerate}
 \item \textbf{边的方向。}有向图用父节点条件分布构造联合分布，无向图用团势构造联合或条件分布。因子图则另外画出因子节点，展示局部计算所需的分解，可以表示这两类模型。
 \item \textbf{节点的取值。}节点可以离散、连续或混合。GMM使用离散成分与连续观测，HMM使用有限状态，线性高斯SSM使用连续状态。节点类型影响求和或积分的计算，但不能单独决定图的方向。
 \item \textbf{条件独立关系。}有向图使用d-分离，无向图使用路径分离。缺边表达在相应条件下被排除的直接依赖，不保证边缘独立。向量节点还需说明内部协方差或分解假设。
 \item \textbf{隐变量安排。}隐变量可能是每个样本的成分、随时间连接的状态，或RBM中的多个二元特征。引入隐变量通常需要边缘化，也可能使观测在边缘上相关。参数不确定性是另一层问题，不应与状态未知混淆。
 \item \textbf{分布参数化。}指数族用固定统计量和自然参数表达概率。高斯局部分布、MEMM的softmax、固定特征CRF和BM联合分布均有相关指数族表达。它不决定有向或无向，也不保证配分函数容易计算。
\end{enumerate}

\begin{table}[htbp]
\centering
\caption{主要图模型的结构、节点与推断对象。}
\label{tab:bayes-model-assumptions}
\small
\begin{tabular}{m{2cm}m{2cm}m{4cm}m{4cm}}
\hline
模型 & 图与节点 & 主要假设 & 参数给定后的计算 \\
\hline
朴素贝叶斯 & 有向星形，节点可混合 & 给定类别后特征独立 & 类别后验的有限求和 \\
GMM & 离散成分、连续观测 & 给定成分后观测高斯 & 成分责任度与混合密度 \\
HMM & 有向链，有限状态 & 一阶状态转移、观测条件独立 & 前向后向与最优路径 \\
线性高斯SSM & 有向链，连续状态 & 线性关系、独立高斯噪声 & 卡尔曼过滤的均值与协方差 \\
MEMM & 条件有向链 & 每个前驱下局部归一化 & 标签路径概率与链式推断 \\
线性链CRF & 条件无向链 & 相邻标签势、全局归一化 & 配分函数、边缘与最优路径 \\
RBM & 无向二部图 & 层内无边，跨层成对能量 & 条件分布精确，联合期望常需采样 \\
\hline
\end{tabular}
\end{table}

表\ref{tab:bayes-model-assumptions}中的模型与推断算法需要分别选择。同一个SSM可以使用不同的近似过滤方法，图结构相同也可以使用不同的局部分布。结构学习、参数学习和变量推断也是不同任务。本章主要假定图结构已选定，再讨论参数学习与变量推断。

\subsection{因子图与消息传递的共同计算方式}
将分布写成$p(x)\propto\prod_a\psi_a(x_a)$后，可以用因子图连接变量节点和因子节点。它表达变量参与哪个局部因子。HMM的条件分布也可以当作因子画入因子图，这不改变原来的联合分布。

\begin{figure}[htbp]
\centering
\begin{tikzpicture}[>=Latex,every node/.style={font=\small},varnode/.style={circle,draw,minimum size=10mm},factor/.style={rectangle,draw,fill=orange!18,minimum size=9mm}]
 \node[varnode] (x1) at (0,0) {$X_1$};
 \node[factor] (f1) at (2,0) {$\psi_1$};
 \node[varnode] (x2) at (4,0) {$X_2$};
 \node[factor] (f2) at (6,0) {$\psi_2$};
 \node[varnode] (x3) at (8,0) {$X_3$};
 \draw[thick] (x1)--(f1)--(x2)--(f2)--(x3);
 \draw[->,blue!65!black] (2.5,0.6)--node[above] {$m_{1\to2}$}(3.5,0.6);
 \draw[->,blue!65!black] (5.5,0.6)--node[above] {$m_{2\to2}$}(4.5,0.6);
 \node at (4,-1) {$p(x_1,x_2,x_3)\propto\psi_1(x_1,x_2)\psi_2(x_2,x_3)$};
\end{tikzpicture}
\caption{树形因子图及向$X_2$汇总的消息。圆形是变量，方形是因子；黑色边只表示变量参与因子，蓝色箭头表示计算消息的方向。两侧分支先对$X_1,X_3$求和，消息相乘并归一化后得到$X_2$的边缘分布。}
\label{fig:bayes-factor-messages}
\end{figure}

图\ref{fig:bayes-factor-messages}中的消息方向是计算安排，不改变原分布的结构。在树形因子图上，边缘化可以利用乘法对加法的分配律逐分支计算。令$N(i)$为变量$i$相邻的因子集合，$N(a)$为因子$a$涉及的变量集合，消息递推为
\begin{align}
 m_{i\to a}(x_i)&=\prod_{b\in N(i)\setminus\{a\}}m_{b\to i}(x_i),\\
 m_{a\to i}(x_i)&=\sum_{x_{N(a)\setminus\{i\}}}\psi_a(x_a)
 \prod_{j\in N(a)\setminus\{i\}}m_{j\to a}(x_j).
\end{align}
消息汇总一个分支对连接变量的权重，不必归一化为概率。叶子向内、再向外计算后，$p(x_i)\propto\prod_{a\in N(i)}m_{a\to i}(x_i)$，最后归一化。连续变量需要积分。HMM和线性链CRF的前向后向就是这种分支消元的具体形式。把求和改为最大值并保存索引，可以恢复最优配置，但这时不再计算边缘概率。

有环图不能直接套用树上的精确性结论。RBM虽然两层条件分布容易计算，跨层连边仍可形成环。一般图精确消元的代价取决于中间因子的大小，常用树宽描述，因此节点数和局部条件分布的简单程度都不足以判断整体推断难度。

\subsection{图模型与贝叶斯参数学习的区别}
在本章的前向后向、卡尔曼递推和RBM条件分布中，参数都暂时视为给定。若还要表达有限数据造成的参数未知，可以在已有图模型外加入参数先验。上一部分线性回归中使用过的贝叶斯公式仍然成立：
\begin{equation}
 p(\theta\mid\mathcal D)\propto p(\mathcal D\mid\theta)p(\theta),\qquad
 p(y_*\mid x_*,\mathcal D)=\int p(y_*\mid x_*,\theta)p(\theta\mid\mathcal D)\,d\theta.
 \label{eq:bayes-predictive-integral}
\end{equation}
第二式假设给定$x_*,\theta$后，新目标与训练数据条件独立，且$x_*$不需要额外更新参数后验。状态求和处理未知状态，参数积分处理未知参数。贝叶斯网络的名称只规定有向分解形式，不要求参数一定采用贝叶斯学习。

\section*{参数不确定性与预测检验}
下面的内容讨论参数后验及预测检验，作为图模型主线之外的补充。首次阅读可先理解前面的图结构与RBM算例，再按需要阅读。这里的参数积分与隐状态边缘化应分别理解。
\paragraph{不确定性的分解与校准。}
在给定模型、输入和数据后，预测方差可用全概率方差公式分解：
\begin{equation}
 \Var(y_*\mid x_*,\mathcal D)=\mathbb E_{p(\theta\mid\mathcal D)}[\Var(y_*\mid x_*,\theta)]
 +\Var_{p(\theta\mid\mathcal D)}[\mathbb E(y_*\mid x_*,\theta)].\label{eq:bayes-predictive-variance}
\end{equation}
式\eqref{eq:bayes-predictive-variance}右端的第一项主要对应观测噪声，第二项对应参数不确定性。模型给出的区间还必须经过覆盖率检查：预测百分之九十区间不应只在训练分布内覆盖百分之六十的真实值。对于工业决策，区间宽度与错误成本同样重要。

\paragraph{全方差公式的推导与估计。}
省略给定的$x_*,\mathcal D$，记$m_\theta=\E[Y_*\mid\theta]$、$m=\E_\theta m_\theta$，则$Y_*-m=(Y_*-m_\theta)+(m_\theta-m)$。平方后取期望，交叉项为零，因为$\E[Y_*-m_\theta\mid\theta]=0$，从而得到全方差公式。这个分解是在指定模型内的，不包含模型类别错设所遗漏的不确定性。

若有$S$个参数样本，并能计算条件均值$m_s$及方差$v_s$，令$\bar m=S^{-1}\sum_sm_s$，则经验混合分布的精确方差为
\begin{equation}
 \hat v=\frac1S\sum_sv_s+\frac1S\sum_s(m_s-\bar m)^2.
\end{equation}
若目的是无偏估计独立后验样本的均值间方差，第二项可改用$S-1$分母，但这时不再是该有限混合的精确方差。回归区间需从完整预测混合求分位数，不能仅对$m_s$取分位数后声称覆盖有噪声观测。对于重尾或多峰预测，$\bar m\pm1.96\sqrt{\hat v}$也不一定是95\%区间。

\paragraph{参数抽样与观测抽样要分两层。}
设两个等权参数样本对应的条件预测分别为
\begin{equation*}
 Y_*\mid\theta^{(1)}\sim\mathcal N(0,1),\qquad
 Y_*\mid\theta^{(2)}\sim\mathcal N(2,1).
\end{equation*}
这里省略共同给定的输入和数据。混合均值为1，平均条件方差为1，条件均值的方差也为1，总预测方差为2。仅保存两个条件均值$(0,2)$会遗漏观测噪声。

生成预测样本时应先选择参数样本，再从对应的条件观测分布抽$Y_*$；生成参数可信区间则只对参数样本取分位数。把抽样次数由100增至10000，只会使这些分位数估计更稳定，并不等于增加了9900次真实设备试验。

\paragraph{贝叶斯方法的常见误用。}
使用贝叶斯模型时，可以从结果倒查计算过程。先检查报告的概率与实际发生频率是否相符，再检查陌生工况下模型是否仍然过度自信。后验方差只反映当前模型表达的未知量，模型没有考虑到的故障机制可能根本不会使它增大。先验也应写明来源和尺度；如果近似算法尚未稳定，计算误差甚至可能大于想要估计的参数不确定性。

\paragraph{预测分布的蒙特卡洛近似。}
当后验积分无法解析计算时，可从后验得到$S$个参数样本$\theta^{(s)}$。使用
\begin{equation}
 p(y_*\mid x_*,\mathcal D)\approx\frac1S\sum_{s=1}^{S}p(y_*\mid x_*,\theta^{(s)}).
\end{equation}
注意先平均概率再取最大类别；平均 logits 或平均类别标签一般不等价。对回归，可由每个参数样本产生一条预测，再计算总方差和分位数区间。$S$太小会造成蒙特卡洛噪声，样本相关时有效样本量还会进一步下降。

\paragraph{采样诊断。}
马尔可夫链蒙特卡洛（Markov Chain Monte Carlo，MCMC）通过具有目标平稳分布的相关样本链近似期望。有效样本量（Effective Sample Size，ESS）衡量某个估计量的采样效率，与粒子滤波中依据归一化权重计算的有效粒子数具有不同定义。
链的轨迹图可以检查是否进入稳定区域，自相关函数衡量样本冗余，有效样本量近似为
\begin{equation}
 ESS\approx\frac{S}{1+2\sum_{k\geq1}\rho_k}.
\end{equation}
多条不同初始状态的链还应比较$\hat R$等收敛指标。若链只在后验的一个模式附近游走，表面上的稳定并不代表全局混合良好。工业高维模型通常更偏向近似推断，但这不意味着可以跳过近似质量和预测覆盖率检查。

\paragraph{Effective Sample Size (ESS)是针对哪个估计量。}
对平稳链上的标量$h(\theta^{(s)})$，设其方差为$v_h$、滞后相关为$\rho_k$，展开样本均值的协方差和得
\begin{equation}
 \Var\left(\frac1S\sum_sh(\theta^{(s)})\right)
 =\frac{v_h}{S}\left[1+2\sum_{k=1}^{S-1}(1-k/S)\rho_k\right].
\end{equation}
当相关可求和且链足够长时，匹配独立样本均值的方差便得到前述ESS。不同参数、尾部概率和维护成本函数有不同的自相关，因此不存在适用于所有结论的单一ESS。短链直接相加所有经验相关会很不稳定，需截断或稳定的谱估计；尚未进入平稳区或未跨模式时，这个局部诊断不能证明准确性。

\paragraph{先验敏感性分析。}
先验选择应通过敏感性分析验证。对故障率、噪声尺度或参数收缩强度使用几组合理先验，比较后验预测和决策是否稳定。若结果对先验极其敏感，说明数据不足、模型可识别性弱，或先验承担了过多结论。此时应增加数据、改进测量或明确报告先验依赖，而不是只选择产生最好结果的先验。

\paragraph{校准、覆盖率与风险。}
预测区间覆盖率和区间宽度需要同时观察。过宽的区间容易达到覆盖率，却可能没有决策价值；过窄的区间看似精确，却会造成过度自信。对分类概率，可按置信度分箱比较预测概率和经验事件率；对回归区间，可按设备、工况和预测提前量分层检查。校准应在独立验证数据上完成，并在部署漂移后持续监测。

\paragraph{小样本故障率的后验决策。}
某设备在新工况下运行$20$次，观察到$2$次异常。若采用$\operatorname{Beta}(1,9)$先验，则后验为$\operatorname{Beta}(3,27)$，后验均值为$0.1$。直接使用样本比例得到的估计也是$0.1$，但二者对未来风险的表达不同：后验还包含小样本造成的宽区间。若把“安排人工复检”的成本记为$C_I$，把一次漏检损失记为$C_M$，则当后验异常概率满足
\begin{equation}
 p(\text{异常}\mid\mathcal D)C_M>C_I
\end{equation}
时才应触发复检。这里假定复检及随后的处理能消除所计漏检损失，且只产生固定成本$C_I$。例如$C_M=100,C_I=12$时，后验概率$0.1$对应不复检损失10，小于12，故不复检；若复检不完美，应重新列出残余损失。这样，概率输出通过成本函数转化为动作，而不是凭经验固定一个看似通用的阈值。

该例还说明先验不是隐藏的常数。若先验改为$\operatorname{Beta}(1,1)$，后验均值变为$3/22$；两种结果的差异提示当前数据不足以完全决定风险。工程报告应同时给出先验、后验和决策阈值，使不同团队可以复核结论。

\paragraph{回归预测区间的覆盖实验。}
假设模型对每台设备给出正态观测预测分布，其均值为$\mu_i$、标准差为$\sigma_i$，则名义95\%预测区间为$[\mu_i-1.96\sigma_i,\mu_i+1.96\sigma_i]$。在验证集上统计真实值落入区间的比例，并分别按设备、环境温度和预测提前量分组。若整体覆盖率为$0.94$，但高温工况只有$0.71$，则总体数字掩盖了局部失效。可以在独立校准集上估计分位数修正因子，再在另一段时间上复验，避免把校准集上的偶然波动误认为改进。

区间宽度也应进入报告。两个模型可能具有相同覆盖率，但一个区间宽得无法指导维护，另一个区间更窄且只在少数工况失效。因而应联合绘制覆盖率、平均区间宽度和漏覆盖成本，并说明校准是否使用了未来数据。

\paragraph{后验预测检查的操作步骤。}
用于拟合后验的真实数据也参与了这项检查，因此它不是独立测试集上的泛化评价。复制数据只是模型内的模拟结果，不能作为新设备证据再次更新同一后验。
后验预测检查不是把训练数据重新拟合一次，而是从后验参数样本生成复制数据$\tilde y$，比较$\tilde y$与真实数据的统计量。例如可比较均值、方差、极值数量、故障持续时间和相邻残差相关性。若模型能匹配均值却无法产生真实的长尾，说明其似然假设不适合风险评估。工业案例中，检查统计量应优先选择与维护动作相关的量，而不是只选择容易展示的总体误差。

本章核心内容可结合下面的图模型课程和RBM实现文档学习。参数后验教程用于补充阅读，不能替代条件独立、分解与隐变量推断的学习。
\section*{辅助学习材料}
\begingroup
\emergencystretch=3em
\begin{itemize}
\item \textbf{Probabilistic Graphical Models}；作者/平台：Stanford University CS228；英文；课程讲义与课程资源。重点学习 Bayesian networks、Markov networks、因子分解和推断主题，帮助理解本章的 HMM、MRF、消息传递与条件独立。\href{https://ermongroup.github.io/cs228/}{课程主页}。
\item \textbf{Bayesian Regression - Introduction (Part 1)}；作者/平台：Pyro 社区；英文；官方教程。阅读模型与先验、变分推断和后验预测部分，将点估计与参数不确定性作对照。\href{https://pyro.ai/examples/bayesian\_regression.html}{在线阅读}。
\item \textbf{BernoulliRBM}；scikit-learn；英文；模型接口文档。对照二元可见层与隐层的定义、参数矩阵形状以及持续对比散度训练，区分隐藏特征的条件均值与随机采样。\href{https://scikit-learn.org/stable/modules/generated/sklearn.neural\_network.BernoulliRBM.html}{接口文档}。
\item \textbf{Deep Learning：概率模型与配分函数相关章节}；Ian Goodfellow、Yoshua Bengio、Aaron Courville；英文；作者公开的在线教材。结合本章RBM自由能与似然梯度，学习能量模型中的归一化和采样计算。\href{https://www.deeplearningbook.org/}{书籍网站}。
\end{itemize}
\endgroup
\paragraph{本章小结。}
概率图模型把变量、条件独立关系与分布分解联系起来。有向图以父节点条件分布相乘，用d-分离判断图保证的独立性。无向图以团势相乘，通过一个配分函数归一化，用路径分离判断条件独立。图规定允许哪些依赖，局部分布和参数决定具体概率。

朴素贝叶斯假设给定类别后特征独立。隐藏类别并采用高斯观测得到GMM，将成分变量沿时间连接得到状态空间模型。有限离散状态可用HMM前向后向推断，线性高斯状态可用卡尔曼滤波，较一般的状态分布可用粒子滤波近似。MEMM直接学习标签条件分布，允许相关输入，但局部归一化会产生标签偏置。CRF使用全局归一化，线性链CRF保留相邻标签的无向马尔可夫关系。

BM在二元无向图中使用成对能量并可包含隐变量，完整变量的联合分布属于指数族。RBM进一步删除层内连边，使两层的条件分布可以分别计算，但边缘依赖、全局配分函数和模型期望仍需处理。边的方向、节点取值、条件独立关系、隐变量安排和指数族参数化，是比较这些模型的五个方面，不能相互替代。

下一章转向参数化与非参数化模型。这里已经用图结构规定多个变量如何共同分布，后面将研究如何直接规定多个输入位置的函数值共同服从高斯分布。高斯条件化仍是计算工具，改变的是被建模的对象与相关结构。
